\documentclass[10pt,a4paper]{amsart}
\usepackage[margin=19mm]{geometry}
\usepackage{amsmath,amssymb,mathtools}
\usepackage[scr=rsfs,cal=euler]{mathalfa}
\usepackage{xfrac}
\usepackage[unicode,hidelinks,bookmarks=false]{hyperref}
\newcommand{\ie}{i.e.\ }
\newcommand{\cf}{cf.\ }
\newcommand{\resp}{respectively\ }
\newcommand{\Subsection}{\subsection}
\providecommand{\nobreakdash}{\nobreak\hbox{-}}
\setlength{\emergencystretch}{2em}
\multlinegap0pt
\usepackage{bm}
\usepackage{bbm}
\usepackage{dsfont}
\usepackage{xspace}
\usepackage{cancel}
\usepackage{mathtools}
\usepackage{cleveref}
\usepackage{amsthm}
\makeatletter
\@ifundefined{lemma}{\newtheorem{lemma}{Lemma}}{}
\@ifundefined{theorem}{\newtheorem{theorem}{Theorem}}{}
\makeatother
\newcommand{\Cycles}{\operatorname{Cyc}}
\newcommand{\Forest}{\operatorname{For}}
\title[2-switch: transition and stability on forests and pseudofest]{2-switch: transition and stability on forests and pseudofests}
\author{Victor N. Schv\"ollner, Adri\'an Pastine, Daniel A. Jaume}
\date{Source: arXiv:2603.07439v3 (CC BY 4.0)}
\begin{document}
\maketitle

\noindent\emph{Source and licence.} 2-switch: transition and stability on forests and pseudofests. Version arXiv:2603.07439v3 (CC BY 4.0), distributed under the Creative Commons Attribution 4.0 International licence, as recorded in the arXiv licence metadata for this exact version and shown on the abstract page https://arxiv.org/abs/2603.07439v3. Full rights notice: licenses/arxiv-2603.07439-CC-BY-4.0.txt.\par\smallskip

\noindent\textit{Editorial scope.} Counted unit: the Lemma whose source label is \texttt{lema2switchpswitch} in the article (source0.tex lines 304--313, source label \texttt{lema2switchpswitch}), delivered together with the article's own complete original proof of it (source0.tex lines 315--332, one \texttt{proof} environment of 18 lines). No mathematics is written, repaired or re-derived: statement and proof are the article's own text line for line. Required context retained uncounted: uswitchcaract (lines 241--251); 2switch.in.C\_U=pswitch.in.U (lines 293--296). Every definition, display and label of those lines is retained unchanged. Omitted from this package and not redistributed: 1--240, 252--292, 297--303, 314--314, 333--1071. Nothing inside a retained excerpt is omitted or altered: every retained body is delivered exactly as the article's lines are written, so no omission receipt of any kind accompanies this package. Citations: the retained excerpts carry no citation command at all, so no bibliography entry of the article is transcribed and no reference is redistributed. Reference adaptation: none was needed and none was made; every reference inside the delivered unit resolves inside it, so no unresolved reference or citation is produced. Printed slips kept literal: no printed word, symbol, hypothesis or number of the counted statement or of its proof was altered. Layout and class: the article's own class is \texttt{BICA} with packages including tikz, float, xcolor, aliascnt, algorithm2e, multicol, cleveref; the wrapper is the lane's pinned \texttt{amsart} editorial wrapper at 10pt on A4, which restates the theorem-like environments the retained text needs and adds the article's own macro definitions that the retained text needs (\texttt{\textbackslash Cycles}, \texttt{\textbackslash Forest}), with extra wrapper packages (bm, bbm, dsfont, xspace, cancel, mathtools, cleveref). The delivered numbering is the unit's own. Assets: no retained span contains a figure environment, a graphics-inclusion command, a PDF-page inclusion, a table or a TikZ picture; no image file is copied into this package, the package declares no asset dependency and no image file is redistributed with it. Licence: version arXiv:2603.07439v3 of this article is distributed under the Creative Commons Attribution 4.0 International licence, as recorded in the arXiv licence metadata for this exact version and shown on the abstract page https://arxiv.org/abs/2603.07439v3. A full rights notice is delivered at licenses/arxiv-2603.07439-CC-BY-4.0.txt. Characters: every retained excerpt is pure ASCII, so the delivered engine, XeLaTeX with its default Latin Modern text font, reports no missing character.
\begin{theorem}
	\label{uswitchcaract}
	Let $\tau$ be a 2-switch between two disjoint edges $e_{1}, e_{2}$ of a unicyclic graph $U$. Then, the following statements hold:
	\begin{enumerate}
		\item If $e_{1},e_{2}\in E(\Forest(U))$, then \(\tau\) is a u-switch on $U$ if and only if $\tau$ is a t-switch on $U-e$, for all $e\in E(\Cycles(U))$.

		\item If $e_{1}\in E(\Cycles(U))$ and $e_{2}\in E(\Forest(U))$, then \(\tau\) is a u-switch on $U$.

		\item If $e_{1},e_{2}\in E(\Cycles(U))$, then \(\tau\) is a u-switch on $U$ if and only if $\tau(\Cycles(U))\approx \Cycles(U)$.
	\end{enumerate}
\end{theorem}

\begin{lemma}
	\label{2switch.in.C_U=pswitch.in.U}
     If $U$ is a unicyclic graph, then every 2-switch between two edges of $\Cycles(U)$ is a p-switch on $U$.
\end{lemma}

\begin{lemma}
	\label{lema2switchpswitch}
	Let $\tau= {{a \ b}\choose{c \ d}}$ be a 2-switch on a pseudoforest $G$ with $c(G)=1$. Suppose that one of the following conditions holds:
	\begin{enumerate}
		\item $ab\in E(\Forest(G))$ and $cd\in E(\Cycles(G))$;

		\item $ab,cd\in E(\Cycles(G))$.
	\end{enumerate}
	Then, $\tau$ is a p-switch on $G$.
\end{lemma}

\begin{proof}
	For each case of the hypothesis we have the following subcases:
	(A) $ab$ and $cd$ lie in the same component of $G$;
	(B) $ab$ and $cd$ lie in different components of $G$.
	\begin{enumerate}
		\item[(1.A)] Use \Cref{uswitchcaract}.

		\item[(2.A)] Use \Cref{2switch.in.C_U=pswitch.in.U}.

		\item[(1.B)] Let $H$ be the component of $G$ containing $ab$, and let $U$ be the component of $G$ containing $cd$. Since $cd \in E(\mathrm{Cyc}(G))$, the component $U$ is unicyclic. Since $ab \in E(\mathrm{For}(G))$, the edge $ab$ does not lie on any cycle of $G$; in particular, $ab$ is a bridge of $H$. By the case hypothesis, $H \ne U$.

		We analyze $\tau(H \,\dot\cup\, U)$ directly. Since $ab$ is a bridge, $H - ab$ has exactly two components: if $H$ is a tree, both are trees; if $H$ is unicyclic with cycle $\mathrm{Cyc}(H)$, one of them contains $\mathrm{Cyc}(H)$ and is therefore unicyclic, while the other is a tree. Let $H_a$ and $H_b$ denote the components of $H - ab$ containing $a$ and $b$, respectively. On the other hand, since $cd \in E(\mathrm{Cyc}(U))$, the graph $T_U = U - cd$ is a tree.

		The 2-switch $\tau$ adds $ac$ and $bd$ to $H_a \,\dot\cup\, H_b \,\dot\cup\, T_U$. Since $a \in V(H_a)$, $b \in V(H_b)$, $c, d \in V(T_U)$, and $V(H_a)$, $V(H_b)$, $V(T_U)$ are pairwise disjoint, the edges $ac$ and $bd$ link the three pieces in a path-like fashion ($H_a$---$T_U$---$H_b$) without creating cycles among them. Hence $\tau(H \,\dot\cup\, U)$ is connected and contains at most one cycle (namely $\mathrm{Cyc}(H)$, if $H$ is unicyclic). The remaining components of $G$ are not affected by $\tau$. Therefore $\tau(G)$ is a pseudoforest.

		\item[(2.B)] Note that $\tau$ glues the two cycles containing $ab$ and $cd$ together into a new cycle. Hence, $c(\tau(G)) = c(G)$. \qedhere
	\end{enumerate}
\end{proof}

\end{document}
